\section{Module \texttt{container.c}}

Le module \texttt{container.c} rassemble les conteneurs de mise en page du wrapper. Il couvre la grille, la boîte linéaire, la pile de pages et deux helpers complémentaires pour les interfaces multi-vues.

\subsection{Fonction \texttt{create\_grid}}

La fonction \texttt{create\_grid} instancie un \texttt{GtkGrid} puis applique l'espacement entre lignes et colonnes ainsi que les options d'homogénéité.

\begin{lstlisting}[language=C]
GtkWidget *create_grid(const grid_config *config) {
    GtkWidget *grid = gtk_grid_new();
    if (config == NULL) {
        return grid;
    }

    gtk_grid_set_column_spacing(GTK_GRID(grid), config->column_spacing);
    gtk_grid_set_row_spacing(GTK_GRID(grid), config->row_spacing);
    gtk_grid_set_column_homogeneous(GTK_GRID(grid), config->homogeneous_columns);
    gtk_grid_set_row_homogeneous(GTK_GRID(grid), config->homogeneous_rows);
    apply_widget_style(grid, &config->style);
    return grid;
}
\end{lstlisting}

\subsection{Fonction \texttt{create\_box}}

\begin{lstlisting}[language=C]
GtkWidget *create_box(const box_config *config) {
    GtkWidget *box = gtk_box_new(
        config ? config->orientation : GTK_ORIENTATION_VERTICAL,
        config ? config->spacing : 0
    );
    if (config == NULL) {
        return box;
    }

    gtk_box_set_homogeneous(GTK_BOX(box), config->homogeneous);
    apply_widget_style(box, &config->style);
    return box;
}
\end{lstlisting}

\subsection{Fonction \texttt{create\_stack}}

\begin{lstlisting}[language=C]
GtkWidget *create_stack(const stack_config *config) {
    GtkWidget *stack = gtk_stack_new();
    if (config == NULL) {
        return stack;
    }

    gtk_stack_set_transition_type(GTK_STACK(stack), config->transition);
    gtk_stack_set_transition_duration(GTK_STACK(stack), config->duration_ms);
    gtk_stack_set_vhomogeneous(GTK_STACK(stack), config->vhomogeneous);
    gtk_stack_set_hhomogeneous(GTK_STACK(stack), config->hhomogeneous);
    apply_widget_style(stack, &config->style);
    return stack;
}
\end{lstlisting}

\subsection{Helpers complémentaires}

Le module ajoute aussi deux helpers pratiques : \texttt{create\_stack\_sidebar}, qui relie une barre latérale à une pile existante, et \texttt{stack\_add\_page}, qui ajoute une page puis renseigne son titre par propriété GTK.

\begin{lstlisting}[language=C]
GtkWidget *create_stack_sidebar(GtkWidget *stack, const widget_style_config *style) {
    GtkWidget *sidebar = gtk_stack_sidebar_new();
    if (stack != NULL) {
        gtk_stack_sidebar_set_stack(GTK_STACK_SIDEBAR(sidebar), GTK_STACK(stack));
    }
    apply_widget_style(sidebar, style);
    return sidebar;
}

void stack_add_page(GtkWidget *stack, GtkWidget *child, const char *name, const char *title) {
    if (stack == NULL || child == NULL || name == NULL || title == NULL) {
        return;
    }

    gtk_stack_add_named(GTK_STACK(stack), child, name);
    g_object_set(gtk_stack_get_page(GTK_STACK(stack), child), "title", title, NULL);
}
\end{lstlisting}

\subsection{APIs GTK utilisées}

\begin{itemize}
    \item \textbf{\texttt{gtk\_grid\_new()}} : Instancie un conteneur en grille à deux dimensions permettant d'organiser les widgets selon des lignes et des colonnes. Elle est le point de départ de la fonction \texttt{create\_grid}, qui configure ensuite l'espacement et l'homogénéité.
    \item \textbf{\texttt{gtk\_box\_new()}} : Crée un conteneur linéaire simple qui dispose ses widgets enfants selon une orientation donnée. Le module l'utilise dans \texttt{create\_box} pour simplifier la mise en place de flux horizontaux ou verticaux.
    \item \textbf{\texttt{gtk\_stack\_new()}} : Prépare un conteneur multicouche affichant un seul widget à la fois parmi plusieurs pages. Cette API est au cœur de \texttt{create\_stack}, facilitant la gestion des transitions et de l'homogénéité des dimensions des pages.
    \item \textbf{\texttt{gtk\_stack\_sidebar\_new()}} : Génère une barre latérale de navigation compatible avec le widget \texttt{GtkStack}. Elle est utilisée par le helper correspondant pour fournir un moyen intuitif de changer de page dans une interface à onglets.
    \item \textbf{\texttt{gtk\_stack\_add\_named()}} : Insère un widget dans une pile en lui assignant un nom symbolique unique. Le wrapper l'utilise dans \texttt{stack\_add\_page} pour permettre un référencement aisé de chaque vue de l'application.
    \item \textbf{\texttt{g\_object\_set()}} : Modifie les propriétés d'un objet via une interface générique à base de chaînes de caractères. Dans le module, elle sert spécifiquement à définir le titre visible d'une page de pile, propriété portée par l'objet \texttt{GtkStackPage} interne à GTK.
\end{itemize}

\newpage
